\chapter{预推免：第二轮机会}
\label{ch08}
\ist{预推免}
\begin{keybox}
预推免是夏令营之后的第二轮机会，也是很多人在9月真正拿到录取结果的阶段。但它的规则和夏令营完全不同：夏令营还有相对统一的信息节奏，预推免则几乎完全由各校自定——有的学校叫它“预推免”，有的叫“预报名”，有的干脆把夏令营和预推免放在同一个服务平台，还有的学校没有独立的预推免环节。所以本章的一条硬纪律是：任何一条“什么时候报名、考什么、什么时候出结果”的结论，都必须点明是哪所学校、哪个学院、哪一年的做法，并且提醒你以该校当年通知为准。

本章覆盖七类问题：投递策略、系统操作、失败应对、录取规则、名额判断、考核形式、承诺风险。全章反复出现的三个时间锚点是：预推免报名多集中在7—9月且各校差异极大；教育部推免服务系统2027年于9月21日9:00开放填报志愿、10月20日截止；系统内的节奏是9月22日8:00起发复试通知、9月24日8:00起发待录取、14:00起确认。九推则指系统开放后的正式报名与录取环节，与预推免首尾相接但不是一回事。

最容易踩的三个坑：把口头承诺当录取（教育部明确“相关承诺一律无效”，一切以系统备案为准）；错过极短的确认时限（有的学院只给1—4小时）；以及“已被推免录取后不得再报名当年硕士统考”这条资格互斥规则。读本章时，建议你先看scope\_note确认适用范围，再看experience\_note决定怎么用。
\end{keybox}

\section{投递策略}
\begin{qa}{M0420}{预推免还需要重新投材料吗，夏令营的材料能改吗}
要重新走一遍投递，但不少学校做了账号与材料承接。\textbf{上海财经大学计算机与人工智能学院2026年预推免通知}写明：申请人在该校夏令营项目里填写的信息与上传材料会自动同步到推免生预报名账户，但提交状态显示为未提交，须复核完善后再点“确认并提交”。\textbf{武汉大学文学院2026年预报名通知}写明，参加过武大2025年暑期夏令营的考生可直接用夏令营账号、密码登录预推免平台，完成材料上传与志愿填报。
所以关键在复核：能同步的不必重做，不能同步的（新成绩单与排名证明、需重新签字盖章的表格、新增获奖与论文）要按新通知补齐。也有学校完全不承接，要求整套重投。以该校当年通知为准。
\nref{S0302}
\nref{S0097}
\end{qa}

\iin{计算机与人工智能学院}
\ith{材料差异}
\ist{预推免}
\begin{qa}{M0428}{预推免要不要再联系一次导师}
建议再联系一次，但要发“更新”而不是“重问”。导师在预推免环节确实可能握有复试资格的门槛：\textbf{西安电子科技大学光电工程学院2027年预推免通知}明确要求考生“联系确认意向导师”，并规定“考生在推免生预报名系统提交材料且导师审核通过后方可进行复试考核”。\textbf{上海财经大学计算机与人工智能学院2026年预推免通知}也把“专家推荐信2封”列为必备材料，说明学院层面仍看重师生匹配。
邮件写三段：你是谁（夏令营或上次联系的时间点）、这几个月新增了什么（最新排名、论文、竞赛、实习）、预推免已报名哪个方向以及需不需要补材料。不要只发“老师我还有机会吗”。
\nref{S0328}
\nref{S0302}
\end{qa}

\iin{西安电子科技大学}
\ith{联系导师}
\ist{预推免}
\begin{qa}{M0447}{预推免要不要降档投学校}
建议用“冲稳保”三档结构，而不是整体降档。预推免录取普遍带滚动性质：\textbf{大连理工大学商学院2027年章程}写明学院“分批次审核‘全国推免服务系统’中申请信息，本着先申请先审核的原则，择优选拔，额满为止”；\textbf{湖南师范大学2027年简章}同样写“分批审核申请信息，实行先申请先审核的原则”。越晚投，可选池越小。
也别把预推免想得太悲观：\textbf{东华大学2026年9月20日推免流程说明}附的实时录取情况里，纺织、材料、化工、机械、环境、物理、信息与智能科学等学院都写着“各专业继续招收，欢迎报考”。
所以边缘生的正确做法是：保底档投够数量，冲档留1—2所，稳档放主力。
\nref{S0550}
\nref{S0605}
\nref{S0419}
\end{qa}

\iin{大连理工大学}
\ith{投递策略}
\ist{预推免}
\begin{qa}{M0448}{夏令营拿了B校优营，预推免还要冲A校吗}
要冲，但先把B校优营的含金量核实清楚。\textbf{上海交通大学人工智能学院2026年夏令营通知}写明：优秀营员仍须“申请我院并符合我校推免生招生条件，学院将根据接收规模从高分到低分择优录取，额满为止”——优营不是录取。\textbf{上海交通大学上海高级金融学院2026年入学推免招生办法}也要求“获得高金优秀营员资格的考生必须按学院规定的具体要求在教育部推免服务系统报名，完成复试及录取信息确认等程序，否则视同放弃”。
所以要确认三件事：B校优营是否绑定专业和学位类型、是否需要在系统开放后极短时间内确认、B校往年对放弃的容忍度。只有B校时限宽松，冲A校才真正安全。
\nref{S0514}
\nref{S0271}
\nref{S0550}
\end{qa}

\iin{上海交通大学}
\ith{投递策略}
\ist{预推免}
\begin{qa}{M0449}{预推免阶段投递数量要不要比夏令营少}
数量可以少而精，但前提是每一所都投得进、接得住。
全国统一层面的限制出现在系统阶段：\textbf{学信网2025年推免服务系统操作流程}规定推免生“可同时在系统填报三个平行志愿（不分主次）”，每个志愿提交后48小时内不允许修改。而预推免是系统开放前的各校自办系统，彼此独立，不受3个限制。
所以要不要少投，取决于材料准备度与面试撞期概率——预推免考核高度集中在9月：\textbf{同济大学环境科学与工程学院2027年复试定在9月11—12日}、\textbf{大连理工大学船舶工程学院面试暂定9月16日}，投太多必然冲突。
建议正式投递控制在4—6所，宁可一所投扎实。
\nref{S0444}
\nref{S0313}
\nref{S0576}
\end{qa}

\iin{教育部学生服务与素质发展中心（学信网/研招网）}
\ith{投递数量}
\ist{预推免}
\begin{qa}{M0450}{预推免被拒很多，是不是定位太高}
先别急着否定自己，连续被拒更可能是批次与名额问题。\textbf{大连理工大学商学院2027年章程}写明学院“分批次审核”申请信息、“先申请先审核”、“额满为止”；\textbf{东华大学2026年9月20日流程说明}附各学院实时录取情况，人文学院、外语学院、数学与统计学院、生物与医学工程学院等已明确“推免名额已满”。拒你的可能是名额，不是你。
\textbf{华中科技大学2027年工作办法}还留了口子：“如果相关院系拟招收推免生计划仍有余额，且有预推免考核合格生源，可按照前述程序继续通过‘推免服务系统’接收考生。”
归因方法：材料没过初审就改材料；复试后被刷就改表达与专业匹配度；根本没进名单才考虑降档。
\nref{S0550}
\nref{S0419}
\nref{S0199}
\end{qa}

\iin{大连理工大学}
\ith{投递策略}
\ist{预推免}
\begin{qa}{M0451}{预推免的录取通知和优营冲突了怎么选}
判断标准只有一个：哪个是“系统里能落地的”，而不是“纸面或口头的”。\textbf{教育部教学厅〔2014〕5号通知}明确“最终推免生录取名单及新生学籍注册均以‘推免服务系统’备案信息为准。未经招生单位公示及‘推免服务系统’备案的考生不得录取”，并且“凡违反政策开展的推荐或录取工作，以及进行的相关承诺，一律无效”。
排优先级：
\begin{enumerate}
  \item 已通过学校考核、系统里能收到待录取通知的 → 优先；
  \item 优营但明确“须在系统报名并完成复试及录取确认程序，否则视同放弃”的 → 视为未定；
  \item 只有口头承诺的 → 不作为保底依据。
\end{enumerate}
另注意\textbf{华南理工大学2027年通知}：“在‘推免服务系统’上与我校达成待录取状态的推免生不得撤销待录取”，一旦确认没有回头路。
\nref{S0048}
\nref{S0479}
\nref{S0271}
\end{qa}

\iin{中华人民共和国教育部}
\ith{投递策略}
\ist{预推免}
\begin{qa}{M0452}{预推免要不要重新找推荐人签字}
看该校清单，不能一概而论；但申请直博基本都要新签。\textbf{上海财经大学计算机与人工智能学院2026年预推免通知}的必备材料含“专家推荐信2封PDF彩色扫描件，需要2位副教授（含）以上老师分别推荐（考生填写报名信息时，在推荐人信息维护界面填写、保存并下载，专家填写签字后扫描）”。\textbf{北京大学环境与能源学院2027年接收推免说明}要求“专家推荐信2封……推荐信须密封并在封口骑缝处签字”。
也有学校预推免不要推荐信：\textbf{天津大学环境科学与工程学院2027年预报名通知}只要求在天津大学研究生招生信息网系统内填报，未列推荐信要求。
做法是拿到目标校材料清单再决定找不找老师。
\nref{S0302}
\nref{S0310}
\nref{S0587}
\end{qa}

\iin{计算机与人工智能学院}
\ith{材料差异}
\ist{预推免}
\begin{qa}{M0461}{预推免要不要交报名费}
分两层看：预推免报名阶段一般不收费；教育部推免服务系统里有一次全国统一的报名费。\textbf{学信网《2027推免生网上支付报名费注意事项》}规定：“推免生需网上支付报名费。推免服务系统提供代收推免报名费服务，报名费标准和规则均以推荐高校所在省级教育招生考试机构规定为准。”\textbf{上海财经大学2027年操作说明}进一步写明：“教育部规定推免生应网上支付报名费，但根据我校安排，本校推免生一律免收相关费用。”
在各校自办的预推免系统里，\textbf{上海财经大学计算机与人工智能学院2026年}、\textbf{天津大学环境科学与工程学院2027年}、\textbf{西安电子科技大学光电工程学院2027年}的预推免通知中均未把报名费列为申请环节。以该校当年通知与推荐高校所在省规定为准。
\nref{S0451}
\nref{S0182}
\nref{S0419}
\nref{S0302}
\nref{S0587}
\nref{S0328}
\end{qa}

\iin{教育部学生服务与素质发展中心（学信网/研招网）}
\ith{报名流程}
\ist{预推免}
\begin{qa}{M0463}{预推免面试撞期了怎么选}
优先保录取确定性更高、承诺更能兑现的那一场，另一场尽早申请改期。预推免复试时间由学院自定，撞期几乎是常态：\textbf{同济大学环境科学与工程学院2027年预推免通知}把复试定在9月11—12日、地点写明明净楼与济阳楼；\textbf{大连理工大学船舶工程学院2027年时间安排}把面试暂定9月16日，并要求加QQ群接收后续安排。
判断顺序：已进入拟录取或递补名单的优先；名额更少、更可能额满为止的优先；两所都未定时，保底档优先于冲档。
改期要在收到复试通知当天就问，越早越容易协调。以该校当年通知为准。
\nref{S0313}
\nref{S0576}
\nref{S0550}
\end{qa}

\iin{同济大学}
\ith{时间冲突}
\ist{预推免}
\begin{qa}{M0465}{预推免阶段还能联系夏令营拒绝过我的导师吗}
可以联系，但只在你有实质新进展时联系。预推免阶段导师对复试资格仍有影响力：\textbf{西安电子科技大学光电工程学院2027年预推免通知}规定，考生“联系确认意向导师”，且“在推免生预报名系统提交材料且导师审核通过后方可进行复试考核”。
分寸：
\begin{enumerate}
  \item 不解释上次为什么被拒，不追问评审细节；
  \item 只报新增内容（排名变化、论文录用、竞赛、新的实习或课程设计）；
  \item 说明这次是通过预推免正式报名，问是否需要补充材料；
  \item 不要求对方立刻答复。
\end{enumerate}
被拒过一次不等于进黑名单，很多课题组夏令营没名额、预推免才有空额。但别在同一轮反复投同一位老师。
\nref{S0328}
\end{qa}

\iin{西安电子科技大学}
\ith{联系导师}
\ist{预推免}
\begin{qa}{M0470}{预推免和本校推免资格公示时间冲突怎么办}
先看规则，再找协调对象。\textbf{教育部教学厅〔2014〕5号通知}规定：“推荐、接收工作在时间上分为互不交叉的两个阶段。推荐工作统一于每年的9月25日前结束，推荐工作结束后启动接收录取工作，接收录取工作统一于每年的10月25日前结束。推荐阶段招生单位不得进行与考生签订接收录取协议等接收阶段工作。”
也就是说本校推荐资格公示应早于接收阶段。但校内公示与预推免考核确实可能撞期：\textbf{安徽师范大学推荐2027届推免生工作方案}安排9月6日学院公示、公示期不少于3天；\textbf{中南民族大学2027年通知}安排9月11—17日学校审定公示；而\textbf{同济大学环境科学与工程学院2027年复试在9月11—12日}、\textbf{大连理工大学船舶工程学院面试暂定9月16日}。
遇到冲突，找本院教务处或学工办申请线上确认或委托办理。
\nref{S0048}
\nref{S0646}
\nref{S0594}
\nref{S0313}
\nref{S0576}
\end{qa}

\iin{中华人民共和国教育部}
\ith{时间冲突}
\ist{预推免}
\section{系统操作}
\begin{qa}{M0432}{推免系统开放当天要做哪些操作}
2027年推免服务系统9月21日9:00开放填报志愿（\textbf{同济大学2027年提示}、\textbf{河海大学2027年通知}、\textbf{上海财经大学2027年操作说明}、\textbf{东华大学2026年流程说明}口径一致）。当天要做：
\begin{enumerate}
  \item 确认个人资料与照片已填妥（照片1KB至100KB、高100—720像素、宽90—480像素且小于高度、jpg）；
  \item 网上缴费；
  \item 填报志愿，最多3个平行志愿，不分先后；
  \item 尽早提交，提交后48小时内不可修改；
  \item 打开系统的邮件与短信提醒。
\end{enumerate}
随后节奏：9月22日8:00起可查看复试通知 → 9月24日8:00起发待录取 → 9月24日14:00起确认待录取。\textbf{中山大学2027年办法}另要求9月18日起先行注册。
以当年教育部与该校通知为准。
\nref{S0489}
\nref{S0176}
\nref{S0182}
\nref{S0419}
\nref{S0444}
\nref{S0167}
\end{qa}

\iin{同济大学}
\ith{系统操作}
\ist{预推免}
\begin{qa}{M0433}{系统里可以填几个志愿，志愿有先后顺序吗}
三个平行志愿，不分先后顺序。\textbf{学信网2025年推免服务系统操作流程}：“推免生（有特殊政策要求的专项计划推免生除外）可同时在系统填报三个平行志愿（不分主次）……每个志愿在提交后的48小时内不允许修改。”\textbf{上海财经大学2027年操作说明}表述为“可同时填报三个平行志愿（不区分先后顺序）”，\textbf{东华大学2026年流程说明}同为“三个平行志愿（不区分先后顺序）”。
所以系统里不存在“第一志愿”这个概念——三个志愿对招生单位是并行可见的。
一个例外：\textbf{上海财经大学2027年操作说明}写明“通过支教团专项计划遴选的推免生只能选择报考本校专业志愿”，专项计划推免生另有约束。以当年系统公告与该校通知为准。
\nref{S0444}
\nref{S0182}
\nref{S0419}
\end{qa}

\iin{教育部学生服务与素质发展中心（学信网/研招网）}
\ith{系统操作}
\ist{预推免}
\begin{qa}{M0434}{填了志愿多久会收到复试通知}
全国没有统一的反馈时限，只规定了招生单位的反馈机制与48小时锁定期。\textbf{学信网2025年推免服务系统操作流程}：“每个志愿在提交后的48小时内不允许修改，期间招生单位下载报名信息并决定是否允许考生参加复试。志愿提交48小时后仍未接到复试通知、复试未通过，或拒绝待录取通知的推免生可继续填报其他志愿。志愿提交48小时内，如招生单位通过系统明确拒绝推免生申请，推免生可立即填报其他志愿。”
实际节奏很快。2027年9月21日9:00起填报志愿；\textbf{东华大学2026年流程说明}与\textbf{同济大学2027年提示}均为9月22日8:00起可查看复试通知；\textbf{河海大学2027年通知}为9月22日上午8点开始发放。
个别学院窗口极短：\textbf{大连理工大学商学院2027年章程}要求递补录取考生“须在系统开通后6小时内申报我学院”。
\nref{S0444}
\nref{S0236}
\nref{S0419}
\nref{S0489}
\nref{S0176}
\nref{S0550}
\end{qa}

\iin{教育部学生服务与素质发展中心（学信网/研招网）}
\ith{系统操作}
\ist{预推免}
\begin{qa}{M0435}{收到复试通知多长时间内要确认}
全国只规定“在招生单位规定时间内”确认，未统一时长。\textbf{学信网2025年推免服务系统操作流程}：“如收到复试通知，应在招生单位规定时间内通过系统回复是否同意参加复试，否则招生单位可取消复试通知。未通过系统确认接受复试通知的考生不能被招生单位录取。”
各校自定时限差异很大：
\par\vspace{2pt}\nobreak\begin{center}\footnotesize
\begin{tabularx}{\textwidth}{>{\raggedright\arraybackslash}X>{\raggedright\arraybackslash}X}
\toprule
\textbf{学校（年份）} & \textbf{确认时限} \\ \midrule
大连理工大学商学院 2027 & 学院发出复试通知后4小时内确认 \\
华南理工大学 2027 & 待录取确认最长不超过24小时 \\
同济大学 2027 & 按学院规定时间，未回复可取消 \\
河海大学 2027 & 按规定时间，否则可取消并补录 \\
\bottomrule
\end{tabularx}
\end{center}\vspace{2pt}\par
结论：把目标学院的确认时限写在备忘录里，并确保该时段能上网。
\nref{S0444}
\nref{S0550}
\nref{S0479}
\nref{S0489}
\nref{S0176}
\end{qa}

\iin{教育部学生服务与素质发展中心（学信网/研招网）}
\ith{系统操作}
\ist{预推免}
\begin{qa}{M0436}{待录取通知多久要确认？不确认会怎样}
不确认的后果有两个：录取不生效，且学校可取消你的待录取通知。\textbf{学信网2025年推免服务系统操作流程}：“推免生应在招生单位规定的时间内通过系统答复是否接受录取通知，否则招生单位可取消待录取通知……推免生未通过系统确认接受‘待录取’通知，则录取不生效。”
时限由各校自定，有的非常短：\textbf{华南理工大学2027年通知}规定“考生必须在学院（系）规定的时间内（最长不超过24小时）通过‘推免服务系统’予以确认，否则视为放弃拟录取资格”；\textbf{大连理工大学商学院2027年章程}更细——“对于9月24日14:00前发出的待录取通知，须在9月24日15:00前确认；对于9月24日14:00后发出的待录取通知，须在待录取发出后1小时内确认”。
2027年全国节奏为9月24日8:00起发待录取、14:00起确认。
\nref{S0444}
\nref{S0479}
\nref{S0550}
\nref{S0489}
\nref{S0176}
\end{qa}

\iin{教育部学生服务与素质发展中心（学信网/研招网）}
\ith{系统操作}
\ist{预推免}
\begin{qa}{M0437}{同时收到两个学校的待录取怎么办}
只能接受一个。\textbf{学信网2025年推免服务系统操作流程}：“待录取通知可拒绝多个，但只能接受一个。如接受，推免生即成为该招生单位的待录取考生。推免生接受招生单位发放的待录取通知后，不能再填报其他志愿，或接受其他复试、待录取通知，其他招生单位也不能再向该推免生发放复试或待录取通知。”\textbf{河海大学2027年通知}复述为“只能确认一个招生单位的待录取通知……一旦确认不得更改”。
两个通知同时到手时的顺序：
\begin{enumerate}
  \item 比对各校确认截止时间，先处理更早截止的；
  \item 决定不去的，主动在系统里点“拒绝”或联系学院说明（\textbf{同济大学2027年提示}给了这条路径）；
  \item 确认前不要失联。
\end{enumerate}
\nref{S0444}
\nref{S0182}
\nref{S0176}
\nref{S0489}
\end{qa}

\iin{教育部学生服务与素质发展中心（学信网/研招网）}
\ith{系统操作}
\ist{预推免}
\begin{qa}{M0439}{学校发的待录取通知能不能拒绝}
能拒绝。\textbf{学信网2025年推免服务系统操作流程}明确“待录取通知可拒绝多个，但只能接受一个”。\textbf{同济大学2027年提示}给了操作路径：“如考生放弃我校拟录取资格，请主动与招生学院联系并提供放弃说明，或在教育部推免服务系统中点击‘拒绝’。”
拒绝前要清楚两件事：
\begin{enumerate}
  \item 你放弃的是该校今年一个名额，学校可据此补录。\textbf{河海大学2027年通知}：“否则学校可取消相关通知，各招生单位可实施补录。”
  \item 拒绝和确认是两回事。一旦点“接受”就不可撤销——\textbf{同济大学2027年提示}：“待录取通知一经确认，不予撤销”；\textbf{华南理工大学2027年通知}：“在‘推免服务系统’上与我校达成待录取状态的推免生不得撤销待录取。”
\end{enumerate}
拒绝请趁早，别拖到名额作废。
\nref{S0444}
\nref{S0489}
\nref{S0479}
\nref{S0176}
\end{qa}

\iin{教育部学生服务与素质发展中心（学信网/研招网）}
\ith{系统操作}
\ist{预推免}
\begin{qa}{M0445}{系统里的候补状态是什么意思}
在教育部推免服务系统的官方流程说明中，只描述了“志愿—复试通知—待录取通知”这条主线，没有“候补”这一官方状态；“候补/递补”是学校自己的说法，通常指考核合格但排名在拟录取名额之外，前面有人放弃时按顺序顶上。
院校实例：
\begin{enumerate}
  \item \textbf{华东师范大学心理与认知科学学院2027年综合考核实施细则}写明“方向内未足额完成，从综合考核合格且未进入拟录取名单的候补”中递补；
  \item \textbf{大连理工大学2027年接收推免生章程}把考核结果分为“待录取、递补录取”，递补录取者“须在系统开通后6小时内申报我学院”，否则需再次参加院系面试选拔。
\end{enumerate}
看到“候补”，第一件事是问清递补顺序和联系电话。
\nref{S0444}
\nref{S0295}
\nref{S0550}\par\textbf{本版实操补充：} 候补代表目前未进入正式接收名单，但院系可能在有人放弃或名额调整时按规则联系递补；它不是录取承诺，也没有全国统一的等待顺序和截止时间。主动问清自己是否有明确位次、递补依据、预计通知截止日，以及是否需要继续参加其他考核。保留其他志愿和备选方案，只有收到正式通知并在系统完成要求后，才把状态视为改变。
\end{qa}

\iin{教育部学生服务与素质发展中心（学信网/研招网）}
\ith{系统操作}
\ist{预推免}
\begin{qa}{M0457}{预推免可以同时参加多个学校的复试吗}
可以。\textbf{教育部教学厅〔2014〕5号通知}明确“‘推免服务系统’保障推免生自主报考，考生在规定时间内可自主多次平行报考多个招生单位及专业”，并规定“所有推免生均享有依据招生政策自主选择报考招生单位和专业的权利”，推荐高校“不得将报考本校作为遴选推免生的条件”。
系统阶段的限制只是同时挂3个：\textbf{学信网2025年推免服务系统操作流程}规定可同时填报三个平行志愿，每个志愿提交后48小时内不允许修改。预推免（系统开放前）由各校自办系统，彼此独立，可同时参加多校考核。\textbf{中国矿业大学2027年复试录取公告}就把“未参加我校预报名而在推免服务系统中报考我校并接到复试通知的考生”列为复试对象。
真正的分水岭是接受待录取之后。
\nref{S0048}
\nref{S0444}
\nref{S0627}
\nref{S0176}
\end{qa}

\iin{中华人民共和国教育部}
\ith{系统操作}
\ist{预推免}
\begin{qa}{M0467}{填系统时第一志愿填错了能改吗}
先纠正前提：系统里没有“第一志愿”，三个是平行志愿、不分先后（\textbf{学信网2025年推免服务系统操作流程}：“可同时在系统填报三个平行志愿（不分主次）”；\textbf{上海财经大学2027年操作说明}：“不区分先后顺序”）。
能不能改：提交后48小时内不能修改。\textbf{学信网帮助中心}：“推免生（有特殊政策要求的专项计划推免生除外）志愿填报确认后48小时内不可以修改。若招生单位在48小时内通过系统拒绝推免生申请，或48小时后推免生仍未接到复试通知、复试未通过，则该志愿自动解锁，推免生可继续填报其他志愿。”
补救路径：等48小时自动解锁，或联系该学院研招办请其在系统里明确拒绝以提前解锁，同时把另外两个志愿用足。
\nref{S0236}
\nref{S0444}
\nref{S0182}
\end{qa}

\iin{教育部学生服务与素质发展中心（学信网/研招网）}
\ith{系统操作}
\ist{预推免}
\begin{qa}{M0468}{系统开放前必须拿到offer吗}
不必须，但没有预录取的人在系统阶段只能拼剩余名额。\textbf{山东大学2027年接收推免生招生录取工作办法}：“未参加我校预报名的推免生，可在‘全国推荐免试攻读研究生（免初试、转段）信息公开暨管理服务系统’开通后，登录系统进行报名。我校各招生单位将根据招生计划情况决定是否发放复试通知和组织考核。”\textbf{华中科技大学2027年工作办法}：“如果相关院系拟招收推免生计划仍有余额，且有预推免考核合格生源，可按照前述程序继续通过‘推免服务系统’接收考生。”
但多数学校先把预推免通过的人接进来。\textbf{同济大学2027年提示}整篇面向“已通过我校各学院考核的推免生”，并要求填报志愿与学院确认的拟录取专业、层次、专项计划一致。建议至少有一个预推免结果保底。
\nref{S0631}
\nref{S0199}
\nref{S0489}
\nref{S0173}
\end{qa}

\iin{山东大学}
\ith{系统操作}
\ist{预推免}
\section{失败应对}
\begin{qa}{M0418}{预推免和夏令营到底有什么区别，是不是补录}
不是补录，是独立且更靠近录取的第二轮选拔。三批次关系：
\par\vspace{2pt}\nobreak\begin{center}\footnotesize
\begin{tabularx}{\textwidth}{>{\raggedright\arraybackslash}X>{\raggedright\arraybackslash}X>{\raggedright\arraybackslash}X}
\toprule
\textbf{批次} & \textbf{时间} & \textbf{结论效力} \\ \midrule
夏令营 & 多在6—7月 & 多为“优秀营员”，不等于录取 \\
预推免 & 报名7—9月、考核9月 & 通过后进系统走复试与待录取 \\
九推 & 系统9月21日—10月20日开放期间 & 正式报名、复试、待录取 \\
\bottomrule
\end{tabularx}
\end{center}\vspace{2pt}\par
\textbf{上海交通大学人工智能学院2026年夏令营通知}写明优营“申请我院并符合我校推免生招生条件，学院将根据接收规模从高分到低分择优录取，额满为止”。\textbf{西安电子科技大学光电工程学院2027年预推免通知}：“参加推免生预报名的学生，其复试成绩以预报名时的考核成绩为准。未进行预报名的学生，须在正式报名后参加复试考核。”
名称也不统一：天津大学环境科学与工程学院2027年叫“预报名”，四川大学把夏令营与预推免放在同一平台。是否有独立预推免各校不同，以该校当年通知为准。
\nref{S0514}
\nref{S0328}
\nref{S0587}
\nref{S0474}
\nref{S0489}
\end{qa}

\iin{上海交通大学}
\ith{阶段区别}
\ist{预推免}
\begin{qa}{M0423}{预推免是不是比夏令营更容易上岸}
不一定更容易，但赔率结构不同：门槛略低、名额更少、时限更紧。
门槛：\textbf{武汉大学文学院2026年预报名通知}的接收条件为“总评成绩名列前茅”，外语要求六级425分及以上（或雅思6.5、TOEFL 90），并明确“所在学校须为国内‘双一流’大学或高水平大学”。
名额：预推免时部分名额已被夏令营占用。\textbf{东华大学2026年9月20日实时公示}显示数学与统计、人文、外语、生物与医学工程等学院已“推免名额已满”，但纺织、材料、化工、机械、环境、物理、信息与智能科学等仍“继续招收”。
时限：2027年系统节奏为9月21日填报 → 9月22日复试通知 → 9月24日待录取，留给决策的只有几天。
所以是否更容易，取决于你投的学院还有没有名额。
\nref{S0097}
\nref{S0419}
\nref{S0489}
\end{qa}

\iin{武汉大学}
\ith{阶段区别}
\ist{预推免}
\begin{qa}{M0424}{预推免阶段学校是不是更缺人，能捡漏吗}
部分学校和专业确实有余额，但窗口极短，且冷热在同一所学校内部就分成两半。
最直接的样本是\textbf{东华大学2026年9月20日推免流程说明}附的“各学院推免生录取情况（实时更新）”：纺织、材料、化工、机械、环境、物理、信息与智能科学、国家卓越工程师学院等写“各专业继续招收，欢迎报考”；旭日经济与管理、服装与艺术设计、数学与统计、人文、外语、生物与医学工程等写“推免名额已满”。
另两条机制：\textbf{华南理工大学2027年通知}写“各院系将视第一批复试接收情况确定是否进行第二批复试”；\textbf{华中科技大学2027年工作办法}写“如果相关院系拟招收推免生计划仍有余额，且有预推免考核合格生源，可按照前述程序继续通过‘推免服务系统’接收考生”。
盯学院官网公告与学院研招办电话，别只看学校总通知。
\nref{S0419}
\nref{S0480}
\nref{S0199}
\nref{S0313}
\end{qa}

\iin{yjszs（高校官网）}
\ith{捡漏}
\ist{预推免}
\begin{qa}{M0430}{预推免和九推是一回事吗}
不是一回事。预推免在系统开放前、由各校自办系统完成考核；九推是系统开放后的正式报名与录取环节，两者流程上首尾相接。
2027年锚点：推免服务系统9月21日9:00开放填报志愿、10月20日截止（\textbf{上海财经大学2027年操作说明}、\textbf{东华大学2026年流程说明}、\textbf{中山大学2027年办法}）。系统内节奏为9月22日8:00起发复试通知、9月24日8:00起发待录取、14:00起确认（\textbf{河海大学2027年通知}、\textbf{同济大学2027年提示}）。
预推免的作用是提前锁定。\textbf{西安电子科技大学光电工程学院2027年预推免通知}：“参加推免生预报名的学生，其复试成绩以预报名时的考核成绩为准。未进行预报名的学生，须在正式报名后参加复试考核。”
九推常用来指系统开放初期及后续补录窗口。以该校当年通知为准。
\nref{S0182}
\nref{S0419}
\nref{S0167}
\nref{S0176}
\nref{S0489}
\nref{S0328}
\end{qa}

\iin{上海财经大学}
\ith{阶段区别}
\ist{预推免}
\begin{qa}{M0438}{接受了待录取还能反悔吗，怎么操作}
确认之后不能撤销，这是全国通行规则。\textbf{学信网2025年推免服务系统操作流程}与\textbf{上海财经大学2027年操作说明}都写明：接受待录取后“不能再填报其他志愿，或接受其他复试、待录取通知”。\textbf{同济大学2027年提示}：“待录取通知一经确认，不予撤销。”\textbf{华南理工大学2027年通知}：“在‘推免服务系统’上与我校达成待录取状态的推免生不得撤销待录取。”
唯一“体面”的时机是在确认之前拒绝。\textbf{同济大学2027年提示}给的路径是“主动与招生学院联系并提供放弃说明，或在教育部推免服务系统中点击‘拒绝’”。
另外，被推免录取后不得再报名参加当年硕士研究生考试招生，否则取消推免录取资格（\textbf{山东大学2027年办法}、\textbf{厦门大学2027年办法}均有）。
\nref{S0444}
\nref{S0182}
\nref{S0489}
\nref{S0479}
\nref{S0631}
\nref{S0532}
\end{qa}

\iin{教育部学生服务与素质发展中心（学信网/研招网）}
\ith{违约}
\ist{预推免}
\begin{qa}{M0441}{预推免被放鸽子了还来得及补录吗}
来得及，但窗口很短，且必须主动争取。补录的逻辑是前面的人不确认、名额空出。
\textbf{河海大学2027年通知}：“推免生应在规定时间内按要求接受并确认待录取通知，否则学校可取消相关通知，各招生单位可实施补录。”\textbf{华中科技大学2027年工作办法}：“如果相关院系拟招收推免生计划仍有余额，且有预推免考核合格生源，可按照前述程序继续通过‘推免服务系统’接收考生。”\textbf{大连理工大学2027年章程}把“待录取、递补录取”明确区分，递补者须在系统开通后6小时内申报该学院。
被放鸽子后的动作：立即联系该学院研招办问余额或递补，同步把志愿转向仍有余额的学校，并联系此前给过意向的其他学院。
\nref{S0176}
\nref{S0199}
\nref{S0550}
\nref{S0627}
\end{qa}

\iin{河海大学}
\ith{失败应对}
\ist{预推免}
\begin{qa}{M0443}{九推捡漏有没有技巧，什么时候刷系统}
关键窗口是系统开放后的头几天，各校节奏高度固定，按时间点刷就够了。
2027年全国节奏（\textbf{同济大学2027年提示}、\textbf{河海大学2027年通知}、\textbf{东华大学2026年流程说明}口径一致）：
\begin{enumerate}
  \item 9月21日9:00起填报志愿；
  \item 9月22日8:00起可查看并确认复试通知；
  \item 9月24日8:00起发放待录取通知；
  \item 9月24日14:00起确认待录取。
\end{enumerate}
要盯的信号有三类：\textbf{东华大学}在流程说明里附各学院“继续招收/名额已满”实时表；\textbf{华南理工大学}视第一批复试接收情况决定是否进行第二批复试；\textbf{大连理工大学}章程规定递补录取考生须在系统开通后6小时内申报该学院。
把目标学院的确认时限抄下来（个别只给1—4小时），9月21—25日保持随时可上网。
\nref{S0489}
\nref{S0176}
\nref{S0419}
\nref{S0480}
\nref{S0550}
\end{qa}

\iin{同济大学}
\ith{捡漏}
\ist{预推免}
\begin{qa}{M0444}{推免系统里没被录取是不是就没学上了}
不是。\textbf{东华大学2026年推免流程说明}明确：“截止规定日期仍未落实接收单位的推免生不再保留推免资格，可参加统考报名。”
但代价要算清：
\begin{enumerate}
  \item 推免资格到期作废，之后只能走全国统考；
  \item 接收录取工作按\textbf{教育部教学厅〔2014〕5号通知}“统一于每年的10月25日前结束”，而统考网上报名也在10月，时间衔接得上，但备考起步明显偏晚；
  \item 已被其他学校拟录取的，不得再报名参加当年硕士研究生考试招生，否则取消推免录取资格（\textbf{山东大学2027年办法}等均有）——“放弃推免去考研”这条路并不通。
\end{enumerate}
结论：9月至10月别把志愿火力集中在一两所学校，至少留一所保底。
\nref{S0419}
\nref{S0048}
\nref{S0631}
\end{qa}

\iin{yjszs（高校官网）}
\ith{失败应对}
\ist{预推免}
\begin{qa}{M0472}{放弃待录取去二战或出国有什么影响}
分两种情况。
确认前放弃：系统允许。\textbf{学信网2025年推免服务系统操作流程}规定“待录取通知可拒绝多个，但只能接受一个”；\textbf{同济大学2027年提示}给的路径是联系招生学院并提供放弃说明，或在系统中点击“拒绝”。代价是影响该校补录（\textbf{河海大学2027年通知}：“各招生单位可实施补录”）。
确认后反悔：\textbf{同济大学2027年提示}“待录取通知一经确认，不予撤销”，\textbf{华南理工大学2027年通知}“达成待录取状态的推免生不得撤销待录取”。且\textbf{山东大学2027年办法}、\textbf{厦门大学2027年办法}都规定已被录取的推免生“不得再报名参加当年硕士研究生考试招生，否则将取消其推免录取资格”。
出国：未见统一规定，须与拟录取学院确认是否影响学籍与诚信记录。
\nref{S0444}
\nref{S0489}
\nref{S0479}
\nref{S0176}
\nref{S0631}
\nref{S0532}
\end{qa}

\iin{教育部学生服务与素质发展中心（学信网/研招网）}
\ith{违约}
\ist{预推免}
\begin{qa}{M0473}{九推结束后还没录取，还能走什么渠道}
最现实的渠道是统考。\textbf{东华大学2026年推免流程说明}明确：“截止规定日期仍未落实接收单位的推免生不再保留推免资格，可参加统考报名。”\textbf{教育部教学厅〔2014〕5号通知}规定接收录取工作“统一于每年的10月25日前结束”，而统考网上报名同月在研招网进行，时间上还能接上，但备考时间已经很紧。
其余渠道：支教团、国优计划等专项计划通常在推免阶段就已确定人选；学校层面的补录只能靠主动打电话问学院研招办是否有余额，且必须在系统关闭前完成。\textbf{中国矿业大学2027年复试录取公告}也只把“未参加预报名而直接报名”的考生纳入正常复试流程。
未见“九推结束后另有统一兜底渠道”的规定。
\nref{S0419}
\nref{S0048}
\nref{S0627}
\end{qa}

\iin{yjszs（高校官网）}
\ith{失败应对}
\ist{预推免}
\section{录取规则}
\begin{qa}{M0419}{预推免一般什么时候开始，什么时候截止}
全国没有统一时间，系统环节有统一锚点：教育部推免服务系统2027年9月21日9:00开放填报志愿、10月20日截止（\textbf{上海财经大学2027年操作说明}、\textbf{中山大学2027年办法}）。\textbf{中山大学}另要求9月18日起完成网上注册。
预推免（系统开放前各校自办环节）差异极大，2026—2027年实例：
\par\vspace{2pt}\nobreak\begin{center}\footnotesize
\begin{tabularx}{\textwidth}{>{\raggedright\arraybackslash}X>{\raggedright\arraybackslash}X}
\toprule
\textbf{学校/学院} & \textbf{报名窗口} \\ \midrule
上海财经大学计算机与人工智能学院 2026 & 7月1日9时至8月20日24时 \\
天津大学环境科学与工程学院 2027 & 7月10日12点至9月15日12点 \\
西安电子科技大学光电工程学院 2027 & 7月31日至9月17日 \\
中国矿业大学 2027 & 8月20日至9月18日17:00 \\
武汉大学文学院 2026 & 9月5日12:00至9月10日9:00 \\
\bottomrule
\end{tabularx}
\end{center}\vspace{2pt}\par
各校是否有独立预推免、什么时候，以该校当年通知为准。
\nref{S0182}
\nref{S0167}
\nref{S0419}
\nref{S0302}
\nref{S0587}
\nref{S0328}
\nref{S0173}
\nref{S0097}
\end{qa}

\iin{上海财经大学}
\ith{预推免时间}
\ist{预推免}
\begin{qa}{M0421}{夏令营没参加，直接冲预推免可行吗}
可行，多数学校预推免不要求参加过夏令营，只要求取得推免资格。\textbf{中国矿业大学2027年预报名通知}的申请条件只有“获得本科毕业学校推荐免试资格的优秀应届本科毕业生”“身心健康状况符合规定”等，未见夏令营前置要求；\textbf{华中科技大学2027年工作办法}也写“所有考生均通过预报名系统提交申请材料”。
更关键的是，有些学校把预推免明确定位成夏令营之外的另一条通道：\textbf{上海交通大学船舶海洋与建筑工程学院2027年预推免办法}申请资格第7条就是“未参加我院同学科2027年优才夏令营”；\textbf{机械与动力工程学院2027年预推免办法}则注明“夏令营已获得同专业同类型五星营员的考生，无需重复报名本次预推免”。
所以没参加夏令营不必焦虑，但要注意这类学校的预推免名额可能已被优营部分占用。
\nref{S0173}
\nref{S0199}
\nref{S0575}
\nref{S0262}
\end{qa}

\iin{中国矿业大学}
\ith{路径选择}
\ist{预推免}
\begin{qa}{M0429}{预推免的录取是不是先到先得}
不完全是“先到先得”，更准确的说法是：分批审核、按考核成绩择优、额满为止，滚动色彩很重。
\textbf{大连理工大学商学院2027年接收推免生章程}：“学院分批次审核‘全国推免服务系统’中申请信息，本着先申请先审核的原则，择优选拔，额满为止。”\textbf{湖南师范大学2027年简章}：“我校分批审核申请信息，实行先申请先审核的原则，择优选拔。”
但决定因素是考核成绩。\textbf{西安电子科技大学光电工程学院2027年预推免通知}：“学院根据复试考核结果，结合师生意愿择优确定拟录取名单。”\textbf{中国药科大学2027年工作办法}要求“合格生源充足专业的差额复试比例一般不低于推免生拟招生计划数的150\%”，说明合格者多于名额。
结论：早投能提高被审核到的概率，录取靠成绩与名额。
\nref{S0550}
\nref{S0605}
\nref{S0328}
\nref{S0418}
\end{qa}

\iin{大连理工大学}
\ith{录取规则}
\ist{预推免}
\begin{qa}{M0455}{预推免的结果一般几天出}
没有统一天数，可以分两段看。
校内预推免环节：报名截止到出结果常在1—2周。\textbf{武汉大学文学院2026年预报名通知}：报名9月10日9:00截止 → 考核拟9月18日 → 9月20日12:00前须完成确认；\textbf{同济大学环境科学与工程学院2027年预报名通知}：报名8月31日17:00截止 → 复试9月11—12日；\textbf{大连理工大学船舶工程学院2027年时间安排}：报名9月7日截止 → 面试暂定9月16日。
系统环节：2027年9月21日9:00填报志愿 → 9月22日8:00起复试通知 → 9月24日8:00起待录取 → 14:00起确认（\textbf{同济大学2027年提示}、\textbf{河海大学2027年通知}），最快3天走完。
所以从复试到结果，快则当天到第二天，慢则一周左右。
\nref{S0097}
\nref{S0313}
\nref{S0576}
\nref{S0489}
\nref{S0176}
\end{qa}

\iin{武汉大学}
\ith{预推免时间}
\ist{预推免}
\begin{qa}{M0458}{预推免的夏令营优秀营员优先意味着什么}
意味着你在排序或复试资格上占先，但不等于自动录取。
\textbf{上海交通大学人工智能学院2026年夏令营通知}：“优秀营员如获得所在学校推荐免试研究生资格，申请我院并符合我校推免生招生条件，学院将根据接收规模从高分到低分择优录取，额满为止。”\textbf{上海交通大学上海高级金融学院2026年入学推免招生办法}：“获得高金优秀营员资格的考生必须按学院规定的具体要求在教育部推免服务系统报名，完成复试及录取信息确认等程序，否则视同放弃。”
另一层含义是免考通道：\textbf{上海交通大学机械与动力工程学院2027年预推免办法}注明“夏令营已获得同专业同类型五星营员的考生，无需重复报名本次预推免”。所以看到“优营优先”，要问清是“免复试”还是“排序靠前”。以该校当年通知为准。
\nref{S0514}
\nref{S0271}
\nref{S0262}
\end{qa}

\iin{上海交通大学}
\ith{录取规则}
\ist{预推免}
\begin{qa}{M0462}{预推免材料截止时间卡得很紧怎么应对}
核心办法是提前备好“常备材料包”，再按校微调。截止时间卡得紧是常态：
\textbf{北京大学环境与能源学院2027年接收推免说明}：纸质材料截止9月8日15:00，“逾期送达不予受理”“打印版信息如与网报信息不一致，将被视为不合格材料”。
\textbf{同济大学环境科学与工程学院2027年预报名通知}：未在8月31日17:00前上传材料，网上报名无效。
\textbf{中国矿业大学2027年预报名通知}：系统8月20日至9月18日17:00。
\textbf{上海财经大学计算机与人工智能学院2026年通知}给了标准清单：申请表、个人陈述、2封专家推荐信、成绩单、前五学期排名证明、英语成绩、身份证、学生证、学籍在线验证报告、诚信承诺书，“严格按照顺序做成一个PDF，并在首页附上目录”。
常备件：盖章成绩单与排名证明、六级成绩单、身份证学生证扫描件、学籍在线验证报告。
\nref{S0310}
\nref{S0313}
\nref{S0173}
\nref{S0302}
\end{qa}

\iin{北京大学环境与能源学院}
\ith{预推免时间}
\ist{预推免}
\begin{qa}{M0469}{预推免拿到offer但排名靠后，会被挤掉吗}
有风险，因为“考核合格”和“被录取”之间还隔着名额与排序。
\textbf{中国药科大学2027年工作办法}要求“合格生源充足专业的差额复试比例一般不低于推免生拟招生计划数的150\%”，合格人数本就多于名额。\textbf{大连理工大学商学院2027年章程}把考核结果分为“待录取、递补录取”，并“根据招生计划、考核结果、合格专业依次录取”；\textbf{华中科技大学2027年工作办法}规定复试考核成绩低于60分（按满分100分折算后）不予录取。
保护自己的办法是尽快确认：\textbf{华南理工大学2027年通知}规定“考生必须在学院（系）规定的时间内（最长不超过24小时）通过‘推免服务系统’予以确认，否则视为放弃拟录取资格”。排名靠后时越早确认越安全。
\nref{S0418}
\nref{S0550}
\nref{S0199}
\nref{S0479}
\end{qa}

\iin{中国药科大学}
\ith{录取规则}
\ist{预推免}
\section{名额判断}
\begin{qa}{M0422}{预推免名额还剩多少？怎么判断}
有的学校会公开实时名额，多数不会。判断靠三条线。
一是学校公开的实时表。\textbf{东华大学2026年9月20日推免流程说明}文末附“各学院推免生录取情况（实时更新）”，逐院写明“各专业继续招收”或“推免名额已满”，这是目前最直接的样本。
二是名额上限。\textbf{中国药科大学2027年工作办法}：“拟招收推免生攻读硕士研究生的数量不超过我校2027年硕士研究生招生计划的50\%；拟招收推免生直接攻读博士研究生的数量原则上不超过2027年博士研究生招生计划总数的10\%。”\textbf{西南财经大学2027年接收工作通知}同样规定“原则上不超过我校当年硕士研究生招生总规模的50\%”。
三是公示名单反推：\textbf{大连理工大学计算机科学与技术学院2027年预推免考核结果公示}、\textbf{中南大学计算机学院2027届拟预推免名单公示}，比对专业目录可估算剩余。
\nref{S0419}
\nref{S0418}
\nref{S0483}
\nref{S0107}
\nref{S0124}
\end{qa}

\iin{yjszs（高校官网）}
\ith{名额判断}
\ist{预推免}
\begin{qa}{M0431}{九推还有机会吗}
有，但机会集中在少数学校和专业，窗口以小时计。
\textbf{东华大学2026年9月20日推免流程说明}附各学院实时录取情况，纺织、材料、化工、机械、环境、物理、信息与智能科学等学院写“各专业继续招收，欢迎报考”，说明系统开放时确实有余额。
制度上的口子也有：\textbf{河海大学2027年通知}写“推免生应在规定时间内按要求接受并确认待录取通知，否则学校可取消相关通知，各招生单位可实施补录”；\textbf{华南理工大学2027年通知}写“各院系将视第一批复试接收情况确定是否进行第二批复试”；\textbf{中国矿业大学2027年复试录取公告}把“未参加我校预报名而在推免服务系统中报考我校并接到复试通知的考生”列为复试对象。
2027年系统9月21日开放、10月20日截止，越往后余额越少。
\nref{S0419}
\nref{S0176}
\nref{S0480}
\nref{S0627}
\nref{S0182}
\end{qa}

\iin{yjszs（高校官网）}
\ith{九推}
\ist{预推免}
\begin{qa}{M0442}{九推阶段哪些学校还有名额}
这个问题没有统一名单，只能给出判断方法。
最直接的样本是\textbf{东华大学2026年9月20日推免流程说明}：同一天同一所学校内部冷热分明——纺织、材料、化工、机械、环境、物理、信息与智能科学等写“各专业继续招收，欢迎报考”；旭日经济与管理、服装与艺术设计、数学与统计、人文、外语、生物与医学工程等写“推免名额已满”。
机制上，\textbf{华中科技大学2027年工作办法}给出补录条件：“如果相关院系拟招收推免生计划仍有余额，且有预推免考核合格生源，可按照前述程序继续通过‘推免服务系统’接收考生。”
可操作判断：优先看往年招不满的学院、非热门专业与新设学院；并注意\textbf{东华大学2026年流程说明}写明“未招满的推免名额将用于统招”。未见普遍适用的“九推余量名单”。
\nref{S0419}
\nref{S0199}
\end{qa}

\iin{yjszs（高校官网）}
\ith{九推}
\ist{预推免}
\begin{qa}{M0456}{预推免阶段本校保底还留着吗}
不一定，有的学校明确不为本院学生开预推免通道，有的则把本校生写进申请条件。
排除型：\textbf{上海财经大学计算机与人工智能学院2026年预推免通知}开头即写“本次预推免选拔不招收计算机与人工智能学院本院学生”；\textbf{上海财经大学经济学院2027年通知}与\textbf{金融学院2027年通知}同为“除金融硕士（国际组织人才培养研究方向）外，本次推免预报名工作仅面向外校和本校外院学生，本院本科生推免接收选拔工作将另行组织”——本校本院保底另有渠道，不能在预推免里顺手做。
包含型：\textbf{天津大学环境科学与工程学院2027年预报名通知}申请条件第1条明确写“（含天津大学本校生）”；\textbf{中国矿业大学2027年复试录取公告}的复试对象也含“校内外推免生”。
以该校当年通知为准。
\nref{S0302}
\nref{S0138}
\nref{S0330}
\nref{S0587}
\nref{S0627}
\end{qa}

\iin{计算机与人工智能学院}
\ith{本校保底}
\ist{预推免}
\begin{qa}{M0459}{预推免名额被本校占满了，外校还有戏吗}
有戏，但要看规则边界。\textbf{教育部教学厅〔2014〕5号通知}规定“所有推免生均享有依据招生政策自主选择报考招生单位和专业的权利，所有推免名额（除有特殊政策要求的专项计划外），均可向其他招生单位推荐”，并明确“2014年起，教育部下达推免名额时不再区分学术学位和专业学位，不再设置留校限额”。
口径要留意：2006年《全国普通高等学校推荐优秀应届本科毕业生免试攻读硕士学位研究生工作管理办法（试行）》第十七条曾规定“设有研究生院的高等学校接收本校推免生的人数，不得超过本校推免生总数的65\%”（特殊类型高校可放宽至75\%）。该条针对接收端，与2014年取消推荐端留校限额不冲突，但两份文件都应核对。
确有只招本校的例外：\textbf{华南理工大学2027年预报名通知}写明“工程硕博士培养改革专项、集成电路工程专业智能EDA专项班及‘1+3’西部计划（含研究生支教团）只接受本校优秀推荐免试生报考”。外校生应避开专项，主攻普通计划。
\nref{S0048}
\nref{S0454}
\nref{S0479}
\end{qa}

\iin{中华人民共和国教育部}
\ith{名额判断}
\ist{预推免}
\begin{qa}{M0460}{有的学校预推免只招本校，怎么判断}
看三处，别靠传言。
第一，申请条件有没有明确接收范围。\textbf{天津大学环境科学与工程学院2027年预报名通知}写“（含天津大学本校生）”；\textbf{中国矿业大学2027年复试录取公告}把复试对象列为“预报名且学院审核通过的校内外推免生”，都不分内外。
第二，有没有排除条款。\textbf{上海财经大学计算机与人工智能学院2026年预推免通知}开头即写“本次预推免选拔不招收计算机与人工智能学院本院学生”——反向排除，外校生反而可以放心报。
第三，看专项计划。\textbf{华南理工大学2027年预报名通知}明确“工程硕博士培养改革专项、集成电路工程专业智能EDA专项班及‘1+3’西部计划（含研究生支教团）只接受本校优秀推荐免试生报考”。
未见统一规定，须逐校读当年通知的申请条件与复试对象两条。
\nref{S0587}
\nref{S0627}
\nref{S0302}
\nref{S0479}
\end{qa}

\iin{天津大学}
\ith{名额判断}
\ist{预推免}
\section{考核形式}
\begin{qa}{M0425}{预推免的考核形式和夏令营一样吗}
不一定相同。夏令营偏学术交流与营员考核，预推免更“复试化”，形式完全由各校各院自定。2026—2027年实例：
\par\vspace{2pt}\nobreak\begin{center}\footnotesize
\begin{tabularx}{\textwidth}{>{\raggedright\arraybackslash}X>{\raggedright\arraybackslash}X}
\toprule
\textbf{学校/学院} & \textbf{预推免考核形式} \\ \midrule
武汉大学文学院 2026 & 现场考核，含背景评估＋综合面试 \\
同济大学环境科学与工程学院 2027 & 面试（9月11—12日，现场） \\
西安电子科技大学光电工程学院 2027 & 统一现场面试，不少于20分钟，全程录音录像 \\
上海大学材料科学与工程学院 2027 & 线上面试，英语口语听力30\%＋专业课70\% \\
同济大学法学院 2027 & 面试与笔试相结合 \\
\bottomrule
\end{tabularx}
\end{center}\vspace{2pt}\par
同样是2027年，同济、西电要求现场，上海大学材料学院则线上。\textbf{河海大学2027年工作办法}还要求“一般每位考生面试时间不少于20分钟”，并实行“三随机”机制。以该校当年通知为准。
\nref{S0097}
\nref{S0313}
\nref{S0328}
\nref{S0256}
\nref{S0242}
\nref{S0178}
\end{qa}

\iin{武汉大学}
\ith{考核形式}
\ist{预推免}
\begin{qa}{M0426}{预推免大多是线上吗，会不会只是走过场}
不是“大多线上”——2027年样本里现场考核反而更常见，也不是走过场。
\textbf{西安电子科技大学光电工程学院2027年通知}：“考核安排：统一采取现场面试”“综合面试时间一般不少于20分钟，全程录音录像，并填写《2027年西安电子科技大学推荐免试研究生复试情况记录表》”。
\textbf{武汉大学文学院2026年预报名通知}：“考核包括背景评估、综合面试等，采取现场考核方式……综合考核成绩为百分制，不及格者不予预接收。”
\textbf{同济大学环境科学与工程学院2027年预报名通知}把复试地点写到了明净楼、济阳楼、工程中心大楼。
线上也有：\textbf{上海大学材料科学与工程学院2027年实施办法}多个学科点写明“复试形式：线上面试”。
考核是打分排序的：\textbf{华东师范大学2027年工作办法}要求“按一定差额比例（一般不低于120\%）确定入围综合考核名单”。以该校当年通知为准。
\nref{S0328}
\nref{S0097}
\nref{S0313}
\nref{S0256}
\nref{S0612}
\end{qa}

\iin{西安电子科技大学}
\ith{考核形式}
\ist{预推免}
\begin{qa}{M0427}{预推免面试比夏令营难还是简单}
难度不一定更高，但更“结果导向”——预推免面试直接对应录取名额，问得更贴专业和个人意图，容错率更低。
\textbf{西安电子科技大学光电工程学院2027年预推免通知}写明考核内容：“加强对考生既往学业、一贯表现、科研能力、综合素质和思想品德等情况的全面考查。考查内容包括英语口语与听力、专业知识与知识结构、科研能力与创新意识、综合素质等”，综合面试一般不少于20分钟。
\textbf{河海大学2027年工作办法}要求“充分保障考生的考核时间，一般每位考生面试时间不少于20分钟”，并建立“三随机”机制：随机确定考生次序、随机确定专家组组成人员、随机抽取考核试题。
所以别按夏令营的“聊天式”准备，要按复试标准准备自我介绍、专业问答与科研追问。
\nref{S0328}
\nref{S0178}
\end{qa}

\iin{西安电子科技大学}
\ith{考核形式}
\ist{预推免}
\begin{qa}{M0453}{预推免线上笔试怎么防作弊，会不会更严}
严格程度由学校自定，线上环节普遍有“全程录音录像”的硬性要求。
\textbf{河海大学2027年工作办法}要求“考核全程全景录音并录像”“考核组成员须现场独立署名评分”“评分记录和考生作答情况任何人不得改动”，并建立“三随机”机制：随机确定考生次序、随机确定专家组组成人员、随机抽取考核试题。
\textbf{中国药科大学2027年工作办法}要求“采取随机选定考生次序、随机确定面试专家组组成人员、随机抽取面试试题等措施”，并规定“同一单位同一学科专业同一批次的复试形式应保持一致”。
\textbf{西安电子科技大学光电工程学院2027年通知}写明综合面试“全程录音录像，并填写《2027年西安电子科技大学推荐免试研究生复试情况记录表》”。
未见统一的线上笔试防作弊专门规定。以该校当年通知为准。
\nref{S0178}
\nref{S0418}
\nref{S0328}
\end{qa}

\iin{河海大学}
\ith{考核形式}
\ist{预推免}
\begin{qa}{M0454}{预推免面试问的问题会不会更深}
会更聚焦你个人做过什么和为什么来，而不是泛知识面。
\textbf{大连理工大学商学院2027年章程}把综合面试流程写得很具体：“（1）英文自我介绍（时长3分钟）；（2）5分钟PPT介绍，介绍内容须包含：专业排名信息；本科课程及成绩介绍；外语能力；参与科研创新类项目、竞赛情况；其他信息；（3）评委与考生问答。”这是典型的预推免面试结构。
\textbf{西安电子科技大学光电工程学院2027年预推免通知}列出的考查内容包括“英语口语与听力、专业知识与知识结构、科研能力与创新意识、综合素质等”。
准备重点：排名背后的原因、项目里你具体负责的部分、读研想做的方向。以该校当年通知为准。
\nref{S0550}
\nref{S0328}
\end{qa}

\iin{大连理工大学}
\ith{面试}
\ist{预推免}
\begin{qa}{M0464}{预推免被问你夏令营为什么没来怎么答}
大方回答，给出真实原因，再把话题拉回现在。这个问题通常不是坑——很多学校确实把预推免定位成夏令营之外的通道：\textbf{上海交通大学船舶海洋与建筑工程学院2027年预推免办法}的申请资格第7条就是“未参加我院同学科2027年优才夏令营”；\textbf{机械与动力工程学院2027年预推免办法}则注明“夏令营已获得同专业同类型五星营员的考生，无需重复报名本次预推免”。没参加夏令营在这类学院是常态。
30秒回答框架：
\begin{enumerate}
  \item 客观原因（信息错过、时间冲突、当时未取得入营资格、方向尚在犹豫）；
  \item 但一直关注，做了哪些准备（读过的论文、做过的项目、选过的课）；
  \item 现在预推免第一时间报名，希望在这里读。
\end{enumerate}
避免贬低夏令营或其他学校，也不要编造被拒的故事。
\nref{S0575}
\nref{S0262}
\end{qa}

\iin{上海交通大学}
\ith{面试}
\ist{预推免}
\section{承诺风险}
\begin{qa}{M0440}{预推免拿了口头offer但系统里没发通知怎么办}
口头承诺不作数，一切以推免服务系统为准。\textbf{教育部教学厅〔2014〕5号通知}明确两点：“凡违反政策开展的推荐或录取工作，以及进行的相关承诺，一律无效”；“最终推免生录取名单及新生学籍注册均以‘推免服务系统’备案信息为准。未经招生单位公示及‘推免服务系统’备案的考生不得录取。”
\textbf{武汉大学文学院2026年预报名通知}印证了预录取也会飞：获得“预报名通过”资格者“仍须在……推免服务系统开通后，通过推免服务系统在规定时间内完成推免接收程序，逾期可能会因全校接收推免生的人数达到预期计划而被拒绝”。
做法：保留其他志愿与其他学校的可能，别回绝一切；系统开放后主动确认对方是否发复试通知；收到后按学院规定时限确认。
\nref{S0048}
\nref{S0097}
\nref{S0479}
\end{qa}

\iin{中华人民共和国教育部}
\ith{承诺风险}
\ist{预推免}
\begin{qa}{M0446}{预推免和考研能不能并行准备}
可以并行，但要提前设止损点，因为两者时间高度重叠且互相消耗。
时间冲突是硬约束：预推免考核集中在9月——\textbf{同济大学环境科学与工程学院2027年复试在9月11—12日}、\textbf{大连理工大学船舶工程学院2027年面试暂定9月16日}、\textbf{西安电子科技大学光电工程学院2027年考核预计9月}；而9—10月正是统考强化期。
更硬的约束是资格互斥：\textbf{山东大学2027年办法}与\textbf{厦门大学2027年办法}都规定“已被我校录取的推免生，不得再报名参加当年硕士研究生考试招生，否则将取消其推免录取资格”。也就是说，一旦推免上岸，统考这条路自动关闭。
建议：9月集中投预推免与复试，统考保持“最低维持”（英语单词＋数学基础题），10月中旬推免系统截止后再决定是否全面转统考。
\nref{S0313}
\nref{S0576}
\nref{S0328}
\nref{S0631}
\nref{S0532}
\end{qa}

\iin{同济大学}
\ith{路径关系}
\ist{预推免}
\begin{qa}{M0466}{预推免的offer能不能保留到系统开放}
可以保留到系统开放，但系统里的规定动作一步不落，否则作废。
\textbf{武汉大学文学院2026年预报名通知}：“获得文学院‘预报名通过’资格的考生，仍须在……推免服务系统开通后，通过推免服务系统在规定时间内完成推免接收程序，逾期可能会因全校接收推免生的人数达到预期计划而被拒绝。”
\textbf{南开大学马克思主义学院接收2027年推免生实施细则}更硬：“已取得我院推免预报名综合考核成绩合格的考生，如未在9月21日9:00-14:00报名我院，视为自动放弃拟录取和递补录取资格。”报名窗口只有5小时。
\textbf{华南理工大学2027年预报名通知}还有一条：“通过复试，获得拟录取资格，但未能获得本科所在学校推荐资格的申请人，拟录取资格无效。”
以该校当年通知为准。
\nref{S0097}
\nref{S0133}
\nref{S0479}
\end{qa}

\iin{武汉大学}
\ith{承诺风险}
\ist{预推免}
\begin{qa}{M0471}{推免系统里的待录取和学校口头承诺不一致怎么办}
以系统里的为准。\textbf{教育部教学厅〔2014〕5号通知}明确“最终推免生录取名单及新生学籍注册均以‘推免服务系统’备案信息为准。未经招生单位公示及‘推免服务系统’备案的考生不得录取”，同文还规定“凡违反政策开展的推荐或录取工作，以及进行的相关承诺，一律无效”。
“不一致”最常见的三种情形：专业、研究方向或学位层次不符——\textbf{同济大学2027年提示}要求“报名志愿须与相关招生学院确认的拟录取专业、研究方向、学历层次（硕士/直博）、专项计划保持一致”；专项计划不符；只有口头承诺而系统里没有通知。
处理原则：不要点确认，先联系学院研招办书面说明并要求更正系统信息。\textbf{学信网2025年推免服务系统操作流程}规定“待录取通知可拒绝多个，但只能接受一个”，确认前拒绝是允许的。
\nref{S0048}
\nref{S0489}
\nref{S0444}
\end{qa}

\iin{中华人民共和国教育部}
\ith{承诺风险}
\ist{预推免}
